\pagenumbering{harfi}
\pagestyle{empty}

% --- Dedication ---
\clearpage
\section*{تقديم به: (اختياري)}
\addcontentsline{toc}{section}{تقديم به}
\vspace{1cm}
در این صفحه دانشجو در صورت تمایل می تواند پروژه خود را به شخص یا اشخاصی تقدیم کند.

% --- Acknowledgments ---
\clearpage
\section*{تشكر و قدرداني: (اختياري)}
\addcontentsline{toc}{section}{تشكر و قدرداني}
\vspace{1cm}
در این قسمت دانشجو می تواند در صورت تمایل از شخص یا اشخاصی که وی را در راستای انجام پروژه تخصصی یاری رسانده اند، تقدیر و تشکر نماید.

% --- Abstract ---
\clearpage
\section*{چکيده}
\addcontentsline{toc}{section}{چکيده}
\vspace{0.5cm}
هر گزارش پروژه یا هر پایان نامه ای با چکیده آغاز می شود كه عبارت است از توضیح راجع به موضوع پروژه، شيوه ها و روش انجام پروژه و نتيجه نهایی. كلمات يا عباراتي كه در اين بخش توضيح داده مي‌شود، بايد كاملاً محوري و مرتبط با موضوع پروژه باشند. در متن چکيده، از ارجاع به منابع و اشاره به جداول و نمودارها اجتناب شود.

\vspace{1cm}
\noindent\textbf{واژه‌هاي كليدي:} تعداد كلمات يا عبارات كليدي حداكثر مي‌تواند پنج كلمه يا عبارت باشد.

% --- Automated TOC, LOF, LOT ---
\clearpage
\pagestyle{plain}
\tableofcontents
\clearpage
\listoffigures
\clearpage
\listoftables

% --- Nomenclature / List of Abbreviations ---
\clearpage
\section*{فهرست علائم اختصاري}
\addcontentsline{toc}{section}{فهرست علائم اختصاري}
\vspace{0.5cm}
\begin{table}[h!]
\centering
\renewcommand{\arraystretch}{1.5}
\begin{tabular}{l@{\hspace{3cm}}r}
\toprule
\textbf{کمیت و واحد} & \textbf{شرح فارسی} \\
\midrule
$v \ (\text{m/s})$ & سرعت \\
$a \ (\text{m/s}^2)$ & شتاب \\
$F \ (\text{N})$ & نیرو \\
\bottomrule
\end{tabular}
\end{table}